\documentclass{article}
\usepackage[UTF8]{ctex}
\usepackage{amsmath,amssymb,amsfonts}
\usepackage{enumitem}
\usepackage{geometry}
\geometry{a4paper, margin=1in}
% === 错题本排版优化（不参与题目标识） ===
\setlength{\parindent}{0pt}
\setlength{\parskip}{0.35em plus 0.10em minus 0.08em}
\setlength{\emergencystretch}{3em}
\setlist{itemsep=0.25em, topsep=0.40em, parsep=0pt}
\allowdisplaybreaks[2]
% === 错题本排版优化结束 ===
\begin{document}
\section*{中国人民大学2025年高等代数真题}
\section*{17．中国人民大学2025年高等代数试卷}
\begin{itemize}

\subsection{题目1}
  \item[1．] （10 分）设 \(V=\operatorname{span}\left\{v_{1}, v_{2}, v_{3}, v_{4}\right\}\) ，请计算 \(V\) 的一组正交基。其中
\[
\left(v_{1}, v_{2}, v_{3}, v_{4}\right)=\left(\begin{array}{llll}
1 & 1 & 1 & 1 \\
1 & 1 & 1 & 0 \\
1 & 0 & 2 & 0 \\
1 & 0 & 2 & 1
\end{array}\right)
\]

\subsection{题目2}
  \item[2．] （20 分）在 \(\mathbb{R}^{3}\) 中计算点 \((4,0,0)^{\mathrm{T}},(0,2,0)^{\mathrm{T}},(1,0,1)^{\mathrm{T}}\) 所在的平面方程．
\end{itemize}

\subsection{题目3}
3．（30 分）用最小二乘法计算距离点（－5，－5），（－4，3），（－3，1），（－2，－3），（－1，4）最近的直线方程．

\subsection{题目4}
4．（30 分）\(l_{1}\) 是经过点 \(a=(1,0,0)^{\mathrm{T}}\) 且平行于向量 \(v=(2,3,4)^{\mathrm{T}}\) 的直线，\(l_{2}\) 是经过 \(b=(0,1,0)^{\mathrm{T}}\) 且平行于向量 \(w=(4,3,2)^{\mathrm{T}}\) 的直线，请计算直线 \(l_{1}\) 与 \(l_{2}\) 上点的最小距离．

\subsection{题目5}
5．（30 分）设 \(\mathcal{A}\) 是欧氏空间 \(V\) 上的一个变换，使得对任意的 \(\alpha, \beta \in V\) ，有 \(\langle\mathcal{A} \alpha, \mathcal{A} \beta\rangle=\langle\alpha, \beta\rangle\) ，其中\(\langle\alpha, \beta\rangle\) 为 \(\alpha\) 与 \(\beta\) 的内积．请证明 \(\mathcal{A}\) 是一个正交变换．

\subsection{题目6}
6．（30 分）找出所有整数 \(x\) 满足

\[
\left\{\begin{array}{l}
5 \mid(x-2) ; \\
7 \mid(x-1) ; \\
11 \mid(x-5) .
\end{array}\right.
\]
\end{document}
